\begin{solapartat}
Tenint en compte que la força que accelera les càrregues és $\vec{F}_m = q\vec{v}\times\vec{B}$ \punts{0,1} i aplicant el producte vectorial, arriben a la conclusió que Q$_1$ és negativa i Q$_2$ es positiva.

Esquema \punts{0,3}

\figura[0.5]{sol_fig1.png}

\punts{0,2} $F = qvB\sin 90°$

\punts{0,4} $v = \dfrac{F}{qB} = \dfrac{1,01\times 10^{-12}}{3,20\times 10^{-19}\times 4,5\times 10^{-1}} = 7,01\times 10^{6}\ m/s$
\end{solapartat}

\begin{solapartat}
\punts{0,2} $\Sigma\vec{F} = m\vec{a} \Rightarrow qvB\sin 90° = m\dfrac{v^2}{r} \Rightarrow r = \dfrac{mv}{qB}$

\punts{0,4}
\begin{align*}
r_1 &= \frac{m_1v}{qB} = \frac{5,32\times 10^{-26}\times 7,01\times 10^{6}}{3,20\times 10^{-19}\times 4,5\times 10^{-1}} = 2,59\ m\\
r_2 &= r_1\frac{m_2}{m_1} = 2,59\,\frac{1,73\times 10^{-25}}{5,32\times 10^{-26}} = 8,42\ m
\end{align*}

\punts{0,2} $\left.\begin{aligned} &v = \omega_2r_2\\ &\omega_2 = 2\pi f_2\end{aligned}\right\} f_2 = \dfrac{v}{2\pi r_2}$

\punts{0,2} $f_2 = \dfrac{7,01\times 10^{6}}{2\pi\times 8,42} = 1,32\times 10^{5}\ Hz$
\end{solapartat}
